% Generated draft: requires mathematical and source review.
\documentclass[11pt]{article}
\usepackage[margin=25mm]{geometry}
\usepackage{fontspec}
\setmainfont{Latin Modern Roman}
\usepackage{amsmath,amssymb,mathtools,unicode-math}
\setmathfont{Latin Modern Math}
\usepackage{enumitem,needspace,xurl,hyperref}
\setlength{\parindent}{0pt}
\setlength{\parskip}{5pt}
\setlength{\emergencystretch}{3em}
\newcommand{\N}{\mathbb{N}}
\newcommand{\Z}{\mathbb{Z}}
\newcommand{\Q}{\mathbb{Q}}
\newcommand{\R}{\mathbb{R}}
\newcommand{\C}{\mathbb{C}}
\newcommand{\F}{\mathbb{F}}
\newcommand{\A}{\mathbb{A}}
\newcommand{\PP}{\mathbb{P}}
\newcommand{\E}{\mathbb{E}}
\newcommand{\eps}{\varepsilon}
\newcommand{\dd}{\,\mathrm{d}}
\begin{document}
\section*{UnsolvedMath ID: 30004418}
\textbf{Source problem code:} catalogue entry
\subsection{Definitions and mathematical statement}
The supplied problem asks for a combinatorial description of the map
\[
Z^{\mathbf 1_{P_n}}:\mathcal P(n)\longrightarrow\mathbb N_0,
\qquad \lambda\longmapsto Z^{\mathbf 1_{P_n}}_\lambda,
\]
where \(\mathbb N_0=\{0,1,2,\ldots\}\). The record does not explicitly define \(P_n\), identify its ambient group, or specify the relevant prime or range of \(n\). Consequently, the following gives a conditional interpretation of the notation, rather than an unqualified completion of the missing hypotheses.

Under the usual symmetric-group interpretation, let \(n\geq 1\) be an integer, let \(p\) be a prime, let \(S_n\) be the symmetric group on \(\{1,\ldots,n\}\), and let \(P_n\leq S_n\) be a Sylow \(p\)-subgroup. Define
\[
\mathcal P(n)=\left\{(\lambda_1,\ldots,\lambda_\ell):\ell\geq 1,\ \lambda_1\geq\cdots\geq\lambda_\ell\geq 1,\ \sum_{i=1}^{\ell}\lambda_i=n\right\}.
\]
For \(\lambda\in\mathcal P(n)\), let \(V^\lambda\) be the irreducible complex representation of \(S_n\) indexed by \(\lambda\), with character \(\chi^\lambda\). Let \(\mathbf 1_{P_n}\) denote the trivial character of \(P_n\). The corresponding Sylow branching coefficient is
\[
Z^{\mathbf 1_{P_n}}_\lambda
=\left\langle\operatorname{Res}^{S_n}_{P_n}\chi^\lambda,\mathbf 1_{P_n}\right\rangle_{P_n}
=\dim_{\C}(V^\lambda)^{P_n},
\]
where
\[
\langle\alpha,\beta\rangle_{P_n}
=\frac{1}{|P_n|}\sum_{g\in P_n}\alpha(g)\overline{\beta(g)},
\qquad
(V^\lambda)^{P_n}=\{v\in V^\lambda:gv=v\text{ for every }g\in P_n\}.
\]
Thus, under this interpretation, the requested task is to describe the full multiplicity map combinatorially: give a combinatorial rule determining \(Z^{\mathbf 1_{P_n}}_\lambda\) for every \(\lambda\in\mathcal P(n)\), for the intended values of \(n\) and the intended Sylow prime. The source does not specify a required form of combinatorial rule. The choice or quantification of \(p\), the intended range of \(n\), and confirmation of this interpretation remain unresolved; no particular prime, extension to \(n=0\), or assertion that all primes are intended is imposed here.
\subsection{Short English statement}
Describe combinatorially the multiplicity of the trivial representation of \(P_n\) in each irreducible representation indexed by a partition of \(n\). Under the conventional symmetric-group interpretation, these are dimensions of \(P_n\)-fixed subspaces. The supplied record leaves the Sylow prime, the range of \(n\), and the explicit ambient-group convention unspecified.
\subsection{Sources}
\begin{itemize}
\item Local UnsolvedMath registry: supplied record with permanent ID \(30004418\), title \emph{Combinatorics of Sylow Branching Coefficients}, exact source problem code \texttt{catalogue entry}, in the Oberwolfach Reports Open Problems collection. The proposed year is recorded as 2020; no individual proposer is supplied.
\item Supplied report citation: \emph{Mini-Workshop: Kronecker, Plethysm, and Sylow Branching Coefficients and their Applications to Complexity Theory} (2021), Oberwolfach Reports, pp.\ 517--568, DOI 10.4171/owr/2020/9. \url{https://doi.org/10.4171/owr/2020/9}. The background also labels this source \emph{Oberwolfach Report 9/2020}.
\item Supplied literature reference: Bessenrodt, Giannelli, and Wildon, ``branching multiplicities.'' \url{https://arxiv.org/abs/1712.02642}.
\item Supplied literature reference: Giannelli, Law, Long, and Vallejo, ``Sylow branching.'' \url{https://doi.org/10.1112/blms.12584}.
\item Supplied literature reference: \emph{Oberwolfach Report 9/2020}. \url{https://doi.org/10.4171/owr/2020/9}.
\item Literature assessment dated exactly \texttt{2026-08-22}: the supplied curation labels the problem \emph{partially solved} and states that induced constituents and several special Sylow-branching families have explicit descriptions, while no general combinatorial formula for the full multiplicity map is known. This is supplied dated curation, described in the background as literature triage derived from the supplied OWR curation; it has not been independently verified here. The source requests independent verification of time-sensitive status and cited claims.
\end{itemize}
\end{document}
